%!TEX root = ../../note.tex
\subsection{The Axioms for the Real Numbers}\label{sec:axioms}

The real numbers form a \term{complete ordered field}. The axioms come in
three groups: the \emph{field axioms} govern $+$ and $\cdot$ (they let us
compute), the \emph{order axioms} govern $<$ (they let us compare), and the
\emph{completeness axiom} is what makes limits possible. The first two
groups hold in $\Q$ as well; only completeness separates $\R$ from $\Q$.

\subsubsection{The field axioms}
\begin{definition}\label{def:field}
    A \term{field} is a set $F$ with two operations
    $+,\,\cdot : F \times F \to F$ such that for all $a,b,c \in F$:
    \begin{enumerate}[label=\textbf{F\arabic*.}, leftmargin=2.5em]
        \item $a + (b + c) = (a + b) + c$;
        \item $a + b = b + a$;
        \item there is $0 \in F$ with $a + 0 = a$;
        \item for each $a$ there is $-a \in F$ with $a + (-a) = 0$;
        \item $a (b c) = (a b) c$;
        \item $a b = b a$;
        \item there is $1 \in F$, $1 \neq 0$, with $a \cdot 1 = a$;
        \item for each $a \neq 0$ there is $a^{-1} \in F$ with $a a^{-1} = 1$;
        \item $a (b + c) = a b + a c$.
    \end{enumerate}
\end{definition}

$\Q$ is a field, so these axioms say nothing about size or order. What they
do is justify every routine algebraic manipulation. Even ``obvious'' facts
must be derived from them:

\begin{theorem}
    Let $a,b,u,z \in \R$.
    \begin{enumerate}[label=(\alph*)]
        \item If $z + a = a$, then $z = 0$.
        \item If $b \neq 0$ and $u b = b$, then $u = 1$.
        \item $a \cdot 0 = 0$.
    \end{enumerate}
\end{theorem}

\begin{proof}
    (a) $z = z + 0 = z + (a - a) = (z + a) - a = a - a = 0$.

    (b) $u = u \cdot 1 = u (b b^{-1}) = (u b) b^{-1} = b b^{-1} = 1$.

    (c) $a\cdot0=a\cdot(0+0)=a\cdot0+a\cdot0$; adding $-(a\cdot0)$ to both
    sides gives $0=a\cdot0$.
\end{proof}

\subsubsection{The order axioms}
\begin{definition}\label{def:order}
    There is a nonempty set $P\subseteq\R$, the \term{positive real
    numbers}, such that
    \begin{enumerate}[label=\textbf{O\arabic*.}, leftmargin=2.5em]
        \item (\emph{trichotomy}) for every $a \in \R$ exactly one of
        $a = 0$, $a \in P$, $-a \in P$ holds;
        \item $a, b \in P \Rightarrow a + b \in P$;
        \item $a, b \in P \Rightarrow a b \in P$.
    \end{enumerate}
    We write $a>b$ (or $b<a$) if $a-b\in P$, and $a\geq b$ (or $b\leq a$)
    if $a-b\in P\cup\set{0}$.
\end{definition}

Every inequality is thus a statement about membership in $P$, and every
manipulation of inequalities reduces to O1--O3. Trichotomy is what makes
case analysis possible, as in the following example.

\begin{example}
    If $a\neq0$, then $a^2>0$.
\end{example}

\begin{proof}
    By trichotomy, $a\in P$ or $-a\in P$. If $a\in P$, then $a^2=a\cdot a\in
    P$ by O3. If $-a\in P$, then $(-a)(-a)\in P$ by O3, and
    $(-a)(-a)=a^2$: multiplying $a+(-a)=0$ by $-a$ gives
    $(-a)a+(-a)(-a)=0$, so $(-a)(-a)=-\bigl((-a)a\bigr)=-(-a^2)=a^2$.
\end{proof}

The positive rationals also satisfy O1--O3, so $\Q$ is an ordered field
too. What it lacks is completeness (\S\ref{sec:completeness}).

\subsubsection{Absolute value}
\begin{definition}
    The \term{absolute value} of $a\in\R$ is
    $\abs{a} = a$ if $a \geq 0$ and $\abs{a} = -a$ if $a < 0$.
\end{definition}

Geometrically, $\abs{a-b}$ is the distance between $a$ and $b$ on the line.

\begin{theorem}\label{thm:abs}
    For all $a, b \in \R$ and $c>0$:
    \begin{enumerate}[label=(\alph*)]
        \item $\abs{ab} = \abs{a}\abs{b}$;
        \item $\abs{a}^2 = a^2$;
        \item $\abs{a} \leq c \iff -c \leq a \leq c$;
        \item $-\abs{a} \leq a \leq \abs{a}$.
    \end{enumerate}
\end{theorem}

\begin{proof}
    Each part follows by splitting into the cases $a\geq0$ and $a<0$.
    Parts (b) and (d) are immediate from the definition.

    (a) If $a,b\geq0$, then $ab\geq0$ and $\abs{ab}=ab=\abs a\abs b$. If
    $a\geq0>b$, then $ab\leq0$ and $\abs{ab}=-ab=a(-b)=\abs a\abs b$;
    the case $b\geq0>a$ is symmetric. If $a,b<0$, then $ab>0$ and
    $\abs{ab}=ab=(-a)(-b)=\abs a\abs b$.

    (c) Let $\abs a\leq c$. If $a\geq0$, then $a=\abs a\leq c$ and
    $-c<0\leq a$. If $a<0$, then $-a=\abs a\leq c$, i.e.\ $-c\leq a$, and
    $a<0<c$. Conversely, let $-c\leq a\leq c$. If $a\geq0$, then
    $\abs a=a\leq c$; if $a<0$, then $-c\leq a$ gives $\abs a=-a\leq c$.
\end{proof}

Part (c) converts one inequality for $\abs{x}$ into two ordinary
inequalities for $x$ and back; it will be used constantly. Its first
application:

\begin{proposition}[Triangle inequality]\label{prop:triangle}
    $\abs{a + b } \leq \abs{a} + \abs{b}$ for all $a, b \in \R$.
\end{proposition}

\begin{proof}
    By Theorem~\ref{thm:abs}(d), $-\abs{a} \leq a \leq \abs{a}$ and
    $-\abs{b} \leq b \leq \abs{b}$. Adding,
    \begin{equation*}
        -\paren{\abs{a} + \abs{b}} \leq a + b \leq \abs{a} + \abs{b}.
    \end{equation*}
    If $\abs a+\abs b=0$, then $a+b=0$ and there is nothing to prove.
    Otherwise Theorem~\ref{thm:abs}(c), with $c=\abs{a}+\abs{b}$, gives
    $\abs{a+b}\leq\abs a+\abs b$.
\end{proof}

\begin{corollary}\label{cor:triangle}
    For all $a,b\in\R$:
    \begin{enumerate}[label=(\arabic*)]
        \item $\abs{\abs{a} - \abs{b}} \leq \abs{a - b}$;
        \item $\abs{a - b} \leq \abs{a} + \abs{b}$.
    \end{enumerate}
\end{corollary}

\begin{proof}
    (1) Apply the triangle inequality to $a=(a-b)+b$ and to $b=(b-a)+a$:
    \begin{equation*}
        \abs a-\abs b\leq\abs{a-b},
        \qquad
        \abs b-\abs a\leq\abs{b-a}=\abs{a-b}.
    \end{equation*}
    Hence $-\abs{a-b}\leq\abs a-\abs b\leq\abs{a-b}$, which is (1) by
    Theorem~\ref{thm:abs}(c).

    (2) $\abs{a - b} = \abs{a + (-b)} \leq \abs{a} + \abs{-b}
    = \abs{a} + \abs{b}$.
\end{proof}

\subsubsection{The completeness axiom}
We state the third axiom now. The notions it uses are made precise in \S\ref{sec:bounds}, and its consequences are the
subject of \S\ref{sec:completeness}.

\begin{axiom}[Completeness]
    Every nonempty subset of $\R$ that is bounded above has a least upper
    bound (supremum) in $\R$.
\end{axiom}

$\Q$ does not have this property: the set $\set{q\in\Q : q^2<2}$ is bounded
above in $\Q$ but has no least upper bound in $\Q$ (see
Proposition~\ref{prop:Q-incomplete}).
